\clearpage\thispagestyle{empty}
\pdfbookmark[0]{Quellen zu den Tafeln für die Störungen der Ceres}{ceresstoerungen180303-sources}
\begin{center}\textbf{Quellen zu den Tafeln für die Störungen der Ceres}\end{center}
\noindent Diese KI-Übertragung in heutiges Deutsch folgt Carl Friedrich Gauß’ Brief an Franz Xaver von Zach vom 26. Dezember 1802, den vier nummerierten Anmerkungen und den fünf Tafeln in \href{https://gdz.sub.uni-goettingen.de/id/PPN236007467}{Gauß’ \textit{Werken}, Band~VI} (Göttingen, 1874), S.~235--243. Der Brief erläutert, wie sich die Störungen der Länge, des Radiusvektors und der Breite durch Jupiter zusammenfassen lassen. Brief und Tafeln behalten ihre jeweilige Nummerierung.
\par\smallskip\noindent Der Brief erschien im März 1803 als Beitrag XXXIII der \textit{Monatlichen Correspondenz}, beginnend auf \href{https://api.digitale-sammlungen.de/iiif/image/v2/bsb10538599_00275/full/1450,/0/default.jpg}{S.~259}. Gauß führt das Verfahren auf \href{https://api.digitale-sammlungen.de/iiif/image/v2/bsb10538599_00276/full/1450,/0/default.jpg}{S.~260} ein, gibt die Formeln auf \href{https://api.digitale-sammlungen.de/iiif/image/v2/bsb10538599_00277/full/1450,/0/default.jpg}{S.~261} an und schließt den Brief auf \href{https://api.digitale-sammlungen.de/iiif/image/v2/bsb10538599_00278/full/1450,/0/default.jpg}{S.~262}. Die Zeitschrift druckt den \href{https://api.digitale-sammlungen.de/iiif/image/v2/bsb10538599_00279/full/1450,/0/default.jpg}{Tafeltitel und die vier Anmerkungen auf S.~263}. Auf \href{https://api.digitale-sammlungen.de/iiif/image/v2/bsb10538599_00280/full/1450,/0/default.jpg}{S.~264 beginnt Tafel I}.
\par\smallskip\noindent Gauß verweist auf seinen \href{https://api.digitale-sammlungen.de/iiif/image/v2/bsb10538598_00512/full/1450,/0/default.jpg}{früheren Brief, Band~VI, S.~498}, und auf die Formeln für die \href{https://api.digitale-sammlungen.de/iiif/image/v2/bsb10538598_00509/full/1450,/0/default.jpg}{Länge} und die \href{https://api.digitale-sammlungen.de/iiif/image/v2/bsb10538598_00510/full/1450,/0/default.jpg}{Breite} im Novemberheft 1802.
\par\smallskip\noindent Übertragung und Satz wurden mit KI erstellt. Die Digitalisate der Zeitschrift stammen von der Bayerischen Staatsbibliothek.
\clearpage
